\documentclass[12pt]{article}

% -------------------------
% PAQUETES NECESARIOS
% -------------------------
\usepackage[utf8]{inputenc}
\usepackage[T1]{fontenc}
\usepackage{titlesec}
\usepackage{fancyhdr}
\usepackage{graphicx}
\usepackage{listings}
\usepackage{xcolor}
\usepackage{geometry}
\usepackage[spanish]{babel}
\usepackage{hyperref}

% -------------------------
% CONFIGURACIÓN DE PÁGINA
% -------------------------
\geometry{margin=2.5cm}

% Fuente tipo Times New Roman equivalente
\usepackage{mathptmx}

% Interlineado sencillo (por defecto en LaTeX)

% -------------------------
% CONFIGURACIÓN DE ENCABEZADO Y PIE
% -------------------------
\pagestyle{fancy}
\fancyhf{}

\fancyhead[L]{Multiprogramación paralela}
\fancyhead[R]{2026/2027}

\fancyfoot[L]{Andrey, Abdellah} % <-- EDITAR
\fancyfoot[R]{\thepage}

\renewcommand{\headrulewidth}{0.4pt}
\renewcommand{\footrulewidth}{0.4pt}

% -------------------------
% INICIO DEL DOCUMENTO
% -------------------------
\begin{document}

% =========================
% PORTADA
% =========================
\thispagestyle{empty}



% Bloque arriba a la derecha
\begin{flushright}
Abdellah Bensaidi, Souri \\
Morales Vicuña, Andrey \\
Asignatura: Multiprogramación paralela \\
Grupo: 1 \\
%  Subgrupo: 1.1 \\

\end{flushright}

\vfill

% Imagen centrada
\begin{center}
\includegraphics[width=0.4\textwidth]{sello.png}
\end{center}

\vfill

% Fecha abajo
\begin{center}
\today
\end{center}

\newpage

\newpage
\tableofcontents
\newpage

\section{Ejercicios}

\begin{enumerate}
    \item \textbf{Detectar los errores de memoria presentes en el código mediante el uso de la herramienta
Valgrind. Mostrar el informe devuelto por Valgrind e indicar cómo se han solucionado los
errores.}


\begin{lstlisting}[breaklines=true]
double *d = (double *) malloc(((n*n-n)/2)*sizeof(int));
PARA QUE NO SE CORTE 

\end{lstlisting}

Para generar el informe hemos utilizado Valgrind, específicamente las siguientes directivas:

\begin{verbatim}
valgrind --tool=memcheck --track-origins=yes --leak-check=yes --log-file=valgrind.txt
\end{verbatim}

Estas directivas permiten buscar fugas de memoria y desbordamientos, además de proporcionar información sobre la línea en la que se produce cada error. Finalmente, se genera un informe con los resultados obtenidos. Véase el informe completo en el fichero \texttt{/doc/img/valgrind.txt} o en el Anexo~\ref{anexo:valgrind}.

\begin{verbatim}
==18066== HEAP SUMMARY:
==18066==     in use at exit: 2,159,600 bytes in 201 blocks
==18066==   total heap usage: 206 allocs, 5 frees, 2,171,792 bytes allocated
==18066==
==18066== 161,600 (1,600 direct, 160,000 indirect) bytes in 100 blocks are definitely lost in loss record 2 of 3
==18066==    at 0x4846828: malloc (in /usr/libexec/valgrind/vgpreload_memcheck-amd64-linux.so)
==18066==    by 0x1098F5: aplicar_mh (mh.c:72)
==18066==    by 0x109339: main (main.c:48)
==18066==
==18066== 1,998,000 bytes in 1 blocks are definitely lost in loss record 3 of 3
==18066==    at 0x4846828: malloc (in /usr/libexec/valgrind/vgpreload_memcheck-amd64-linux.so)
==18066==    by 0x109561: read_distances (io.c:14)
==18066==    by 0x1092DE: main (main.c:34)
==18066==
==18066== LEAK SUMMARY:
==18066==    definitely lost: 1,999,600 bytes in 101 blocks
==18066==    indirectly lost: 160,000 bytes in 100 blocks
==18066==      possibly lost: 0 bytes in 0 blocks
==18066==      still reachable: 0 bytes in 0 blocks
==18066==      suppressed: 0 bytes in 0 blocks
==18066==
==18066== For lists of detected and suppressed errors, rerun with: -s
==18066== ERROR SUMMARY: 63941 errors from 34 contexts (suppressed: 0 from 0)
\end{verbatim}

Como podemos ver en el resumen, destaca el uso de la memoria heap: se realizaron 206 asignaciones de memoria y únicamente 5 liberaciones. Además, se detectaron 63941 errores.

A continuación, iremos mostrando cómo identificamos los problemas de memoria, así como las decisiones que tomamos para corregirlos.

\begin{verbatim}
==18066== Invalid write of size 8
==18066==    at 0x48DDC8B: __vfscanf_internal (vfscanf-internal.c:2530)
==18066==    by 0x48CDEC5: __isoc99_scanf (isoc99_scanf.c:30)
==18066==    by 0x1095B2: read_distances (io.c:20)
==18066==    by 0x1092DE: main (main.c:34)
==18066==  Address 0x4c6bcf0 is 0 bytes after a block of size 1,998,000 alloc'd
==18066==    at 0x4846828: malloc (in /usr/libexec/valgrind/vgpreload_memcheck-amd64-linux.so)
\end{verbatim}

El primer error corresponde al fichero \texttt{io.c}, en la línea 14:

\begin{verbatim}
double *d = (double *) malloc(((n*n-n)/2)*sizeof(int));
\end{verbatim}

El problema que tenemos aquí es que se declara un array de tipo \texttt{double} y se reserva memoria utilizando el tamaño de un entero. Por lo tanto, la cantidad de memoria reservada no es la adecuada para los elementos que se almacenarán en el array. La solución es sencilla: cambiar \texttt{int} por \texttt{double}, que es el tipo que corresponde.

\begin{verbatim}
==24819== Invalid read of size 4
==24819==    at 0x109720: printf (stdio2.h:86)
==24819==    by 0x109720: print_solution (io.c:46)
==24819==    by 0x1093A6: main (main.c:59)
==24819==  Address 0x4e54a20 is 0 bytes inside a block of size 1,600 free'd
==24819==    at 0x484988F: free (in /usr/libexec/valgrind/vgpreload_memcheck-amd64-linux.so)
==24819==    by 0x109393: main (main.c:56)
==24819==  Block was alloc'd at
==24819==    at 0x4846828: malloc (in /usr/libexec/valgrind/vgpreload_memcheck-amd64-linux.so)
==24819==    by 0x1092ED: main (main.c:41)
\end{verbatim}

En la línea 56 de \texttt{main.c} se libera la variable mediante \texttt{free(sol);}, pero poco después se vuelve a utilizar esta misma variable al pasarla como parámetro a una función. Para solucionarlo, únicamente tenemos que mover la liberación a una línea posterior, en la que ya no necesitemos utilizarla.

\begin{verbatim}
==28993== Invalid read of size 4
==28993==    at 0x109720: printf (stdio2.h:86)
==28993==    by 0x109720: print_solution (io.c:46)
==28993==    by 0x10939E: main (main.c:58)
==28993==  Address 0x4e55060 is 0 bytes after a block of size 1,600 alloc'd
==28993==    at 0x4846828: malloc (in /usr/libexec/valgrind/vgpreload_memcheck-amd64-linux.so)
==28993==    by 0x1092ED: main (main.c:41)
\end{verbatim}

En la línea 46 de \texttt{io.c} nos salimos de la memoria al utilizar un menor o igual:

\begin{verbatim}
for(int i = 0; i < m; i++) { printf("%d ", solucion[i]); }
\end{verbatim}

Se soluciona quitando el igual para evitar salirnos de los límites del array.

\begin{verbatim}
==31682==  Uninitialised value was created by a heap allocation
==31682==    at 0x4846828: malloc (in /usr/libexec/valgrind/vgpreload_memcheck-amd64-linux.so)
==31682==    by 0x1098F5: aplicar_mh (mh.c:72)
==31682==    by 0x109339: main (main.c:48)
\end{verbatim}

En la línea 122 de \texttt{mh.c} se libera \texttt{poblacion}, pero \texttt{poblacion} contiene punteros a listas, por lo que no estamos liberando todo su contenido. Por este motivo aparecen numerosos errores relacionados con accesos a memoria no inicializada.

Se soluciona liberando todo el contenido antes de liberar \texttt{poblacion}:

\begin{verbatim}
for (int  i=0; i<tam_pob ; i++){
free(poblacion[i]->array_int);// segundo arreglo leak memoria
free(poblacion[i]); // primer arreglo leak memoria. El tercero esta en io.c free d
}
\end{verbatim}

Finalmente, se crea una variable \texttt{d} que nunca se libera. Buscamos el punto en el que deja de utilizarse y la liberamos mediante \texttt{free(d)} al final de \texttt{main}, junto con la liberación de \texttt{sol}.

\begin{verbatim}
==36402== HEAP SUMMARY:
==36402==     in use at exit: 0 bytes in 0 blocks
==36402==   total heap usage: 206 allocs, 206 frees, 4,169,792 bytes allocated
==36402==
==36402== All heap blocks were freed -- no leaks are possible
==36402==
\end{verbatim}

Este último resultado nos confirma que se ha liberado toda la memoria correctamente y que no quedan fugas de memoria.


    \item \textbf{Implementar las funciones cruzar(), mutar() y fitness() contenidas en el fichero
mh.c. Los cambios realizados en el código deben ir acompañados de comentarios. En la
memoria se explicará cómo se ha implementado cada una de las funciones, justificando las
decisiones tomadas. El código mostrado en la memoria debe coincidir con el que se entregue.}

    \item \textbf{Introduce un nuevo parámetro que represente la tasa de mutación, m\_rate, que será leído del
fichero de entrada in.txt, y que sustituirá a la constante definida para tal propósito en el
fichero mh.c, facilitándose así la experimentación con este parámetro que pasa a ser variable.
Habrá que modificar también los ficheros run.sh y Makefile así como los argumentos de
algunas funciones del programa. Entonces, en el fichero de entrada habrá cinco columnas: n,
m, n\_gen, tam\_pob, m\_rate codificando, por ejemplo, los valores 500 200 100 50 0,15.}

    \item \textbf{Cambiar el criterio de convergencia del algoritmo evolutivo de forma que se realice en función
del fitness. El algoritmo finalizará cuando el valor de fitness no mejore el obtenido en la
iteración anterior (o las 5, o 10, etc., iteraciones anteriores si el mejor fitness no cambia
demasiado entre iteraciones) en un valor porcentual establecido. Prueba con valores como
2\%, 5\%, 10\% u otros que estimes oportunos. Los resultados obtenidos con cada uno deben
reflejar tanto el tiempo empleado en obtener la solución como valor final de fitness alcanzado.}


    %\begin{itemize}
    %    \item \textbf{}
    %    \item \textbf{}
    %\end{itemize}

    %\begin{verbatim} 
    %\end{verbatim}


\end{enumerate}



\appendix

\section{Informe completo de Valgrind}
\label{anexo:valgrind}

    \begin{verbatim} 
    ==18066== Memcheck, a memory error detector
==18066== Copyright (C) 2002-2022, and GNU GPL'd, by Julian Seward et al.
==18066== Using Valgrind-3.22.0 and LibVEX; rerun with -h for copyright info
==18066== Command: ./sec 1000 400 200 100
==18066== Parent PID: 18065
==18066== 
==18066== Invalid write of size 8
==18066==    at 0x48DDC8B: __vfscanf_internal (vfscanf-internal.c:2530)
==18066==    by 0x48CDEC5: __isoc99_scanf (isoc99_scanf.c:30)
==18066==    by 0x1095B2: read_distances (io.c:20)
==18066==    by 0x1092DE: main (main.c:34)
==18066==  Address 0x4c6bcf0 is 0 bytes after a block of size 1,998,000 alloc'd
==18066==    at 0x4846828: malloc (in /usr/libexec/valgrind/vgpreload_memcheck-amd64-linux.so)
==18066==    by 0x109561: read_distances (io.c:14)
==18066==    by 0x1092DE: main (main.c:34)
==18066== 
==18066== Conditional jump or move depends on uninitialised value(s)
==18066==    at 0x48B7952: msort_with_tmp.part.0 (qsort.c:239)
==18066==    by 0x48B7740: msort_with_tmp (qsort.c:201)
==18066==    by 0x48B7740: msort_with_tmp.part.0 (qsort.c:210)
==18066==    by 0x48B7724: msort_with_tmp (qsort.c:201)
==18066==    by 0x48B7724: msort_with_tmp.part.0 (qsort.c:209)
==18066==    by 0x48B7724: msort_with_tmp (qsort.c:201)
==18066==    by 0x48B7724: msort_with_tmp.part.0 (qsort.c:209)
==18066==    by 0x48B7724: msort_with_tmp (qsort.c:201)
==18066==    by 0x48B7724: msort_with_tmp.part.0 (qsort.c:209)
==18066==    by 0x48B7724: msort_with_tmp (qsort.c:201)
==18066==    by 0x48B7724: msort_with_tmp.part.0 (qsort.c:209)
==18066==    by 0x48B7724: msort_with_tmp (qsort.c:201)
==18066==    by 0x48B7724: msort_with_tmp.part.0 (qsort.c:209)
==18066==    by 0x48B7C1C: msort_with_tmp (qsort.c:201)
==18066==    by 0x48B7C1C: qsort_r (qsort.c:393)
==18066==    by 0x10992C: aplicar_mh (mh.c:80)
==18066==    by 0x109339: main (main.c:48)
==18066==  Uninitialised value was created by a heap allocation
==18066==    at 0x4846828: malloc (in /usr/libexec/valgrind/vgpreload_memcheck-amd64-linux.so)
==18066==    by 0x1098F5: aplicar_mh (mh.c:72)
==18066==    by 0x109339: main (main.c:48)
==18066== 
==18066== Conditional jump or move depends on uninitialised value(s)
==18066==    at 0x48B7952: msort_with_tmp.part.0 (qsort.c:239)
==18066==    by 0x48B7724: msort_with_tmp (qsort.c:201)
==18066==    by 0x48B7724: msort_with_tmp.part.0 (qsort.c:209)
==18066==    by 0x48B7724: msort_with_tmp (qsort.c:201)
==18066==    by 0x48B7724: msort_with_tmp.part.0 (qsort.c:209)
==18066==    by 0x48B7724: msort_with_tmp (qsort.c:201)
==18066==    by 0x48B7724: msort_with_tmp.part.0 (qsort.c:209)
==18066==    by 0x48B7724: msort_with_tmp (qsort.c:201)
==18066==    by 0x48B7724: msort_with_tmp.part.0 (qsort.c:209)
==18066==    by 0x48B7724: msort_with_tmp (qsort.c:201)
==18066==    by 0x48B7724: msort_with_tmp.part.0 (qsort.c:209)
==18066==    by 0x48B7C1C: msort_with_tmp (qsort.c:201)
==18066==    by 0x48B7C1C: qsort_r (qsort.c:393)
==18066==    by 0x10992C: aplicar_mh (mh.c:80)
==18066==    by 0x109339: main (main.c:48)
==18066==  Uninitialised value was created by a heap allocation
==18066==    at 0x4846828: malloc (in /usr/libexec/valgrind/vgpreload_memcheck-amd64-linux.so)
==18066==    by 0x1098F5: aplicar_mh (mh.c:72)
==18066==    by 0x109339: main (main.c:48)
==18066== 
==18066== Conditional jump or move depends on uninitialised value(s)
==18066==    at 0x48B7952: msort_with_tmp.part.0 (qsort.c:239)
==18066==    by 0x48B7740: msort_with_tmp (qsort.c:201)
==18066==    by 0x48B7740: msort_with_tmp.part.0 (qsort.c:210)
==18066==    by 0x48B7740: msort_with_tmp (qsort.c:201)
==18066==    by 0x48B7740: msort_with_tmp.part.0 (qsort.c:210)
==18066==    by 0x48B7724: msort_with_tmp (qsort.c:201)
==18066==    by 0x48B7724: msort_with_tmp.part.0 (qsort.c:209)
==18066==    by 0x48B7724: msort_with_tmp (qsort.c:201)
==18066==    by 0x48B7724: msort_with_tmp.part.0 (qsort.c:209)
==18066==    by 0x48B7724: msort_with_tmp (qsort.c:201)
==18066==    by 0x48B7724: msort_with_tmp.part.0 (qsort.c:209)
==18066==    by 0x48B7724: msort_with_tmp (qsort.c:201)
==18066==    by 0x48B7724: msort_with_tmp.part.0 (qsort.c:209)
==18066==    by 0x48B7C1C: msort_with_tmp (qsort.c:201)
==18066==    by 0x48B7C1C: qsort_r (qsort.c:393)
==18066==    by 0x10992C: aplicar_mh (mh.c:80)
==18066==    by 0x109339: main (main.c:48)
==18066==  Uninitialised value was created by a heap allocation
==18066==    at 0x4846828: malloc (in /usr/libexec/valgrind/vgpreload_memcheck-amd64-linux.so)
==18066==    by 0x1098F5: aplicar_mh (mh.c:72)
==18066==    by 0x109339: main (main.c:48)
==18066== 
==18066== Conditional jump or move depends on uninitialised value(s)
==18066==    at 0x48B7952: msort_with_tmp.part.0 (qsort.c:239)
==18066==    by 0x48B7740: msort_with_tmp (qsort.c:201)
==18066==    by 0x48B7740: msort_with_tmp.part.0 (qsort.c:210)
==18066==    by 0x48B7724: msort_with_tmp (qsort.c:201)
==18066==    by 0x48B7724: msort_with_tmp.part.0 (qsort.c:209)
==18066==    by 0x48B7740: msort_with_tmp (qsort.c:201)
==18066==    by 0x48B7740: msort_with_tmp.part.0 (qsort.c:210)
==18066==    by 0x48B7724: msort_with_tmp (qsort.c:201)
==18066==    by 0x48B7724: msort_with_tmp.part.0 (qsort.c:209)
==18066==    by 0x48B7724: msort_with_tmp (qsort.c:201)
==18066==    by 0x48B7724: msort_with_tmp.part.0 (qsort.c:209)
==18066==    by 0x48B7724: msort_with_tmp (qsort.c:201)
==18066==    by 0x48B7724: msort_with_tmp.part.0 (qsort.c:209)
==18066==    by 0x48B7C1C: msort_with_tmp (qsort.c:201)
==18066==    by 0x48B7C1C: qsort_r (qsort.c:393)
==18066==    by 0x10992C: aplicar_mh (mh.c:80)
==18066==    by 0x109339: main (main.c:48)
==18066==  Uninitialised value was created by a heap allocation
==18066==    at 0x4846828: malloc (in /usr/libexec/valgrind/vgpreload_memcheck-amd64-linux.so)
==18066==    by 0x1098F5: aplicar_mh (mh.c:72)
==18066==    by 0x109339: main (main.c:48)
==18066== 
==18066== Conditional jump or move depends on uninitialised value(s)
==18066==    at 0x48B7952: msort_with_tmp.part.0 (qsort.c:239)
==18066==    by 0x48B7724: msort_with_tmp (qsort.c:201)
==18066==    by 0x48B7724: msort_with_tmp.part.0 (qsort.c:209)
==18066==    by 0x48B7740: msort_with_tmp (qsort.c:201)
==18066==    by 0x48B7740: msort_with_tmp.part.0 (qsort.c:210)
==18066==    by 0x48B7724: msort_with_tmp (qsort.c:201)
==18066==    by 0x48B7724: msort_with_tmp.part.0 (qsort.c:209)
==18066==    by 0x48B7724: msort_with_tmp (qsort.c:201)
==18066==    by 0x48B7724: msort_with_tmp.part.0 (qsort.c:209)
==18066==    by 0x48B7724: msort_with_tmp (qsort.c:201)
==18066==    by 0x48B7724: msort_with_tmp.part.0 (qsort.c:209)
==18066==    by 0x48B7C1C: msort_with_tmp (qsort.c:201)
==18066==    by 0x48B7C1C: qsort_r (qsort.c:393)
==18066==    by 0x10992C: aplicar_mh (mh.c:80)
==18066==    by 0x109339: main (main.c:48)
==18066==  Uninitialised value was created by a heap allocation
==18066==    at 0x4846828: malloc (in /usr/libexec/valgrind/vgpreload_memcheck-amd64-linux.so)
==18066==    by 0x1098F5: aplicar_mh (mh.c:72)
==18066==    by 0x109339: main (main.c:48)
==18066== 
==18066== Conditional jump or move depends on uninitialised value(s)
==18066==    at 0x48B7952: msort_with_tmp.part.0 (qsort.c:239)
==18066==    by 0x48B7740: msort_with_tmp (qsort.c:201)
==18066==    by 0x48B7740: msort_with_tmp.part.0 (qsort.c:210)
==18066==    by 0x48B7740: msort_with_tmp (qsort.c:201)
==18066==    by 0x48B7740: msort_with_tmp.part.0 (qsort.c:210)
==18066==    by 0x48B7740: msort_with_tmp (qsort.c:201)
==18066==    by 0x48B7740: msort_with_tmp.part.0 (qsort.c:210)
==18066==    by 0x48B7724: msort_with_tmp (qsort.c:201)
==18066==    by 0x48B7724: msort_with_tmp.part.0 (qsort.c:209)
==18066==    by 0x48B7724: msort_with_tmp (qsort.c:201)
==18066==    by 0x48B7724: msort_with_tmp.part.0 (qsort.c:209)
==18066==    by 0x48B7724: msort_with_tmp (qsort.c:201)
==18066==    by 0x48B7724: msort_with_tmp.part.0 (qsort.c:209)
==18066==    by 0x48B7C1C: msort_with_tmp (qsort.c:201)
==18066==    by 0x48B7C1C: qsort_r (qsort.c:393)
==18066==    by 0x10992C: aplicar_mh (mh.c:80)
==18066==    by 0x109339: main (main.c:48)
==18066==  Uninitialised value was created by a heap allocation
==18066==    at 0x4846828: malloc (in /usr/libexec/valgrind/vgpreload_memcheck-amd64-linux.so)
==18066==    by 0x1098F5: aplicar_mh (mh.c:72)
==18066==    by 0x109339: main (main.c:48)
==18066== 
==18066== Conditional jump or move depends on uninitialised value(s)
==18066==    at 0x48B7952: msort_with_tmp.part.0 (qsort.c:239)
==18066==    by 0x48B7724: msort_with_tmp (qsort.c:201)
==18066==    by 0x48B7724: msort_with_tmp.part.0 (qsort.c:209)
==18066==    by 0x48B7724: msort_with_tmp (qsort.c:201)
==18066==    by 0x48B7724: msort_with_tmp.part.0 (qsort.c:209)
==18066==    by 0x48B7740: msort_with_tmp (qsort.c:201)
==18066==    by 0x48B7740: msort_with_tmp.part.0 (qsort.c:210)
==18066==    by 0x48B7724: msort_with_tmp (qsort.c:201)
==18066==    by 0x48B7724: msort_with_tmp.part.0 (qsort.c:209)
==18066==    by 0x48B7724: msort_with_tmp (qsort.c:201)
==18066==    by 0x48B7724: msort_with_tmp.part.0 (qsort.c:209)
==18066==    by 0x48B7C1C: msort_with_tmp (qsort.c:201)
==18066==    by 0x48B7C1C: qsort_r (qsort.c:393)
==18066==    by 0x10992C: aplicar_mh (mh.c:80)
==18066==    by 0x109339: main (main.c:48)
==18066==  Uninitialised value was created by a heap allocation
==18066==    at 0x4846828: malloc (in /usr/libexec/valgrind/vgpreload_memcheck-amd64-linux.so)
==18066==    by 0x1098F5: aplicar_mh (mh.c:72)
==18066==    by 0x109339: main (main.c:48)
==18066== 
==18066== Conditional jump or move depends on uninitialised value(s)
==18066==    at 0x48B7952: msort_with_tmp.part.0 (qsort.c:239)
==18066==    by 0x48B7724: msort_with_tmp (qsort.c:201)
==18066==    by 0x48B7724: msort_with_tmp.part.0 (qsort.c:209)
==18066==    by 0x48B7740: msort_with_tmp (qsort.c:201)
==18066==    by 0x48B7740: msort_with_tmp.part.0 (qsort.c:210)
==18066==    by 0x48B7740: msort_with_tmp (qsort.c:201)
==18066==    by 0x48B7740: msort_with_tmp.part.0 (qsort.c:210)
==18066==    by 0x48B7724: msort_with_tmp (qsort.c:201)
==18066==    by 0x48B7724: msort_with_tmp.part.0 (qsort.c:209)
==18066==    by 0x48B7724: msort_with_tmp (qsort.c:201)
==18066==    by 0x48B7724: msort_with_tmp.part.0 (qsort.c:209)
==18066==    by 0x48B7C1C: msort_with_tmp (qsort.c:201)
==18066==    by 0x48B7C1C: qsort_r (qsort.c:393)
==18066==    by 0x10992C: aplicar_mh (mh.c:80)
==18066==    by 0x109339: main (main.c:48)
==18066==  Uninitialised value was created by a heap allocation
==18066==    at 0x4846828: malloc (in /usr/libexec/valgrind/vgpreload_memcheck-amd64-linux.so)
==18066==    by 0x1098F5: aplicar_mh (mh.c:72)
==18066==    by 0x109339: main (main.c:48)
==18066== 
==18066== Conditional jump or move depends on uninitialised value(s)
==18066==    at 0x48B7952: msort_with_tmp.part.0 (qsort.c:239)
==18066==    by 0x48B7724: msort_with_tmp (qsort.c:201)
==18066==    by 0x48B7724: msort_with_tmp.part.0 (qsort.c:209)
==18066==    by 0x48B7724: msort_with_tmp (qsort.c:201)
==18066==    by 0x48B7724: msort_with_tmp.part.0 (qsort.c:209)
==18066==    by 0x48B7C1C: msort_with_tmp (qsort.c:201)
==18066==    by 0x48B7C1C: qsort_r (qsort.c:393)
==18066==    by 0x10992C: aplicar_mh (mh.c:80)
==18066==    by 0x109339: main (main.c:48)
==18066==  Uninitialised value was created by a heap allocation
==18066==    at 0x4846828: malloc (in /usr/libexec/valgrind/vgpreload_memcheck-amd64-linux.so)
==18066==    by 0x1098F5: aplicar_mh (mh.c:72)
==18066==    by 0x109339: main (main.c:48)
==18066== 
==18066== Conditional jump or move depends on uninitialised value(s)
==18066==    at 0x48B7952: msort_with_tmp.part.0 (qsort.c:239)
==18066==    by 0x48B7740: msort_with_tmp (qsort.c:201)
==18066==    by 0x48B7740: msort_with_tmp.part.0 (qsort.c:210)
==18066==    by 0x48B7724: msort_with_tmp (qsort.c:201)
==18066==    by 0x48B7724: msort_with_tmp.part.0 (qsort.c:209)
==18066==    by 0x48B7C1C: msort_with_tmp (qsort.c:201)
==18066==    by 0x48B7C1C: qsort_r (qsort.c:393)
==18066==    by 0x10992C: aplicar_mh (mh.c:80)
==18066==    by 0x109339: main (main.c:48)
==18066==  Uninitialised value was created by a heap allocation
==18066==    at 0x4846828: malloc (in /usr/libexec/valgrind/vgpreload_memcheck-amd64-linux.so)
==18066==    by 0x1098F5: aplicar_mh (mh.c:72)
==18066==    by 0x109339: main (main.c:48)
==18066== 
==18066== Conditional jump or move depends on uninitialised value(s)
==18066==    at 0x48B7952: msort_with_tmp.part.0 (qsort.c:239)
==18066==    by 0x48B7724: msort_with_tmp (qsort.c:201)
==18066==    by 0x48B7724: msort_with_tmp.part.0 (qsort.c:209)
==18066==    by 0x48B7C1C: msort_with_tmp (qsort.c:201)
==18066==    by 0x48B7C1C: qsort_r (qsort.c:393)
==18066==    by 0x10992C: aplicar_mh (mh.c:80)
==18066==    by 0x109339: main (main.c:48)
==18066==  Uninitialised value was created by a heap allocation
==18066==    at 0x4846828: malloc (in /usr/libexec/valgrind/vgpreload_memcheck-amd64-linux.so)
==18066==    by 0x1098F5: aplicar_mh (mh.c:72)
==18066==    by 0x109339: main (main.c:48)
==18066== 
==18066== Conditional jump or move depends on uninitialised value(s)
==18066==    at 0x48B7952: msort_with_tmp.part.0 (qsort.c:239)
==18066==    by 0x48B7724: msort_with_tmp (qsort.c:201)
==18066==    by 0x48B7724: msort_with_tmp.part.0 (qsort.c:209)
==18066==    by 0x48B7740: msort_with_tmp (qsort.c:201)
==18066==    by 0x48B7740: msort_with_tmp.part.0 (qsort.c:210)
==18066==    by 0x48B7C1C: msort_with_tmp (qsort.c:201)
==18066==    by 0x48B7C1C: qsort_r (qsort.c:393)
==18066==    by 0x10992C: aplicar_mh (mh.c:80)
==18066==    by 0x109339: main (main.c:48)
==18066==  Uninitialised value was created by a heap allocation
==18066==    at 0x4846828: malloc (in /usr/libexec/valgrind/vgpreload_memcheck-amd64-linux.so)
==18066==    by 0x1098F5: aplicar_mh (mh.c:72)
==18066==    by 0x109339: main (main.c:48)
==18066== 
==18066== Conditional jump or move depends on uninitialised value(s)
==18066==    at 0x48B7952: msort_with_tmp.part.0 (qsort.c:239)
==18066==    by 0x48B7740: msort_with_tmp (qsort.c:201)
==18066==    by 0x48B7740: msort_with_tmp.part.0 (qsort.c:210)
==18066==    by 0x48B7740: msort_with_tmp (qsort.c:201)
==18066==    by 0x48B7740: msort_with_tmp.part.0 (qsort.c:210)
==18066==    by 0x48B7C1C: msort_with_tmp (qsort.c:201)
==18066==    by 0x48B7C1C: qsort_r (qsort.c:393)
==18066==    by 0x10992C: aplicar_mh (mh.c:80)
==18066==    by 0x109339: main (main.c:48)
==18066==  Uninitialised value was created by a heap allocation
==18066==    at 0x4846828: malloc (in /usr/libexec/valgrind/vgpreload_memcheck-amd64-linux.so)
==18066==    by 0x1098F5: aplicar_mh (mh.c:72)
==18066==    by 0x109339: main (main.c:48)
==18066== 
==18066== Conditional jump or move depends on uninitialised value(s)
==18066==    at 0x48B7952: msort_with_tmp.part.0 (qsort.c:239)
==18066==    by 0x48B7740: msort_with_tmp (qsort.c:201)
==18066==    by 0x48B7740: msort_with_tmp.part.0 (qsort.c:210)
==18066==    by 0x48B7C1C: msort_with_tmp (qsort.c:201)
==18066==    by 0x48B7C1C: qsort_r (qsort.c:393)
==18066==    by 0x10992C: aplicar_mh (mh.c:80)
==18066==    by 0x109339: main (main.c:48)
==18066==  Uninitialised value was created by a heap allocation
==18066==    at 0x4846828: malloc (in /usr/libexec/valgrind/vgpreload_memcheck-amd64-linux.so)
==18066==    by 0x1098F5: aplicar_mh (mh.c:72)
==18066==    by 0x109339: main (main.c:48)
==18066== 
==18066== Conditional jump or move depends on uninitialised value(s)
==18066==    at 0x48B7952: msort_with_tmp.part.0 (qsort.c:239)
==18066==    by 0x48B7C1C: msort_with_tmp (qsort.c:201)
==18066==    by 0x48B7C1C: qsort_r (qsort.c:393)
==18066==    by 0x10992C: aplicar_mh (mh.c:80)
==18066==    by 0x109339: main (main.c:48)
==18066==  Uninitialised value was created by a heap allocation
==18066==    at 0x4846828: malloc (in /usr/libexec/valgrind/vgpreload_memcheck-amd64-linux.so)
==18066==    by 0x1098F5: aplicar_mh (mh.c:72)
==18066==    by 0x109339: main (main.c:48)
==18066== 
==18066== Conditional jump or move depends on uninitialised value(s)
==18066==    at 0x48B7952: msort_with_tmp.part.0 (qsort.c:239)
==18066==    by 0x48B7724: msort_with_tmp (qsort.c:201)
==18066==    by 0x48B7724: msort_with_tmp.part.0 (qsort.c:209)
==18066==    by 0x48B7C1C: msort_with_tmp (qsort.c:201)
==18066==    by 0x48B7C1C: qsort_r (qsort.c:393)
==18066==    by 0x10994A: aplicar_mh (mh.c:104)
==18066==    by 0x109339: main (main.c:48)
==18066==  Uninitialised value was created by a heap allocation
==18066==    at 0x4846828: malloc (in /usr/libexec/valgrind/vgpreload_memcheck-amd64-linux.so)
==18066==    by 0x1098F5: aplicar_mh (mh.c:72)
==18066==    by 0x109339: main (main.c:48)
==18066== 
==18066== Conditional jump or move depends on uninitialised value(s)
==18066==    at 0x48B7952: msort_with_tmp.part.0 (qsort.c:239)
==18066==    by 0x48B7740: msort_with_tmp (qsort.c:201)
==18066==    by 0x48B7740: msort_with_tmp.part.0 (qsort.c:210)
==18066==    by 0x48B7C1C: msort_with_tmp (qsort.c:201)
==18066==    by 0x48B7C1C: qsort_r (qsort.c:393)
==18066==    by 0x10994A: aplicar_mh (mh.c:104)
==18066==    by 0x109339: main (main.c:48)
==18066==  Uninitialised value was created by a heap allocation
==18066==    at 0x4846828: malloc (in /usr/libexec/valgrind/vgpreload_memcheck-amd64-linux.so)
==18066==    by 0x1098F5: aplicar_mh (mh.c:72)
==18066==    by 0x109339: main (main.c:48)
==18066== 
==18066== Conditional jump or move depends on uninitialised value(s)
==18066==    at 0x48B7952: msort_with_tmp.part.0 (qsort.c:239)
==18066==    by 0x48B7C1C: msort_with_tmp (qsort.c:201)
==18066==    by 0x48B7C1C: qsort_r (qsort.c:393)
==18066==    by 0x10994A: aplicar_mh (mh.c:104)
==18066==    by 0x109339: main (main.c:48)
==18066==  Uninitialised value was created by a heap allocation
==18066==    at 0x4846828: malloc (in /usr/libexec/valgrind/vgpreload_memcheck-amd64-linux.so)
==18066==    by 0x1098F5: aplicar_mh (mh.c:72)
==18066==    by 0x109339: main (main.c:48)
==18066== 
==18066== Invalid read of size 4
==18066==    at 0x109720: printf (stdio2.h:86)
==18066==    by 0x109720: print_solution (io.c:46)
==18066==    by 0x1093A6: main (main.c:59)
==18066==  Address 0x4c6cd70 is 0 bytes inside a block of size 1,600 free'd
==18066==    at 0x484988F: free (in /usr/libexec/valgrind/vgpreload_memcheck-amd64-linux.so)
==18066==    by 0x109393: main (main.c:56)
==18066==  Block was alloc'd at
==18066==    at 0x4846828: malloc (in /usr/libexec/valgrind/vgpreload_memcheck-amd64-linux.so)
==18066==    by 0x1092ED: main (main.c:41)
==18066== 
==18066== Conditional jump or move depends on uninitialised value(s)
==18066==    at 0x48CF181: __printf_fp_buffer_1.isra.0 (printf_fp.c:230)
==18066==    by 0x48D135B: __printf_fp_l_buffer (printf_fp.c:1122)
==18066==    by 0x48D86C4: __printf_fp_spec (vfprintf-internal.c:266)
==18066==    by 0x48D86C4: __printf_buffer (vfprintf-internal.c:999)
==18066==    by 0x48D973A: __vfprintf_internal (vfprintf-internal.c:1544)
==18066==    by 0x49A5C9C: __printf_chk (printf_chk.c:33)
==18066==    by 0x1093A6: main (main.c:59)
==18066==  Uninitialised value was created by a heap allocation
==18066==    at 0x4846828: malloc (in /usr/libexec/valgrind/vgpreload_memcheck-amd64-linux.so)
==18066==    by 0x1098F5: aplicar_mh (mh.c:72)
==18066==    by 0x109339: main (main.c:48)
==18066== 
==18066== Conditional jump or move depends on uninitialised value(s)
==18066==    at 0x48CF19B: __printf_fp_buffer_1.isra.0 (printf_fp.c:230)
==18066==    by 0x48D135B: __printf_fp_l_buffer (printf_fp.c:1122)
==18066==    by 0x48D86C4: __printf_fp_spec (vfprintf-internal.c:266)
==18066==    by 0x48D86C4: __printf_buffer (vfprintf-internal.c:999)
==18066==    by 0x48D973A: __vfprintf_internal (vfprintf-internal.c:1544)
==18066==    by 0x49A5C9C: __printf_chk (printf_chk.c:33)
==18066==    by 0x1093A6: main (main.c:59)
==18066==  Uninitialised value was created by a heap allocation
==18066==    at 0x4846828: malloc (in /usr/libexec/valgrind/vgpreload_memcheck-amd64-linux.so)
==18066==    by 0x1098F5: aplicar_mh (mh.c:72)
==18066==    by 0x109339: main (main.c:48)
==18066== 
==18066== Conditional jump or move depends on uninitialised value(s)
==18066==    at 0x48C8D46: __mpn_extract_double (dbl2mpn.c:56)
==18066==    by 0x48CF925: __printf_fp_buffer_1.isra.0 (printf_fp.c:230)
==18066==    by 0x48D135B: __printf_fp_l_buffer (printf_fp.c:1122)
==18066==    by 0x48D86C4: __printf_fp_spec (vfprintf-internal.c:266)
==18066==    by 0x48D86C4: __printf_buffer (vfprintf-internal.c:999)
==18066==    by 0x48D973A: __vfprintf_internal (vfprintf-internal.c:1544)
==18066==    by 0x49A5C9C: __printf_chk (printf_chk.c:33)
==18066==    by 0x1093A6: main (main.c:59)
==18066==  Uninitialised value was created by a heap allocation
==18066==    at 0x4846828: malloc (in /usr/libexec/valgrind/vgpreload_memcheck-amd64-linux.so)
==18066==    by 0x1098F5: aplicar_mh (mh.c:72)
==18066==    by 0x109339: main (main.c:48)
==18066== 
==18066== Conditional jump or move depends on uninitialised value(s)
==18066==    at 0x48C8D4B: __mpn_extract_double (dbl2mpn.c:60)
==18066==    by 0x48CF925: __printf_fp_buffer_1.isra.0 (printf_fp.c:230)
==18066==    by 0x48D135B: __printf_fp_l_buffer (printf_fp.c:1122)
==18066==    by 0x48D86C4: __printf_fp_spec (vfprintf-internal.c:266)
==18066==    by 0x48D86C4: __printf_buffer (vfprintf-internal.c:999)
==18066==    by 0x48D973A: __vfprintf_internal (vfprintf-internal.c:1544)
==18066==    by 0x49A5C9C: __printf_chk (printf_chk.c:33)
==18066==    by 0x1093A6: main (main.c:59)
==18066==  Uninitialised value was created by a heap allocation
==18066==    at 0x4846828: malloc (in /usr/libexec/valgrind/vgpreload_memcheck-amd64-linux.so)
==18066==    by 0x1098F5: aplicar_mh (mh.c:72)
==18066==    by 0x109339: main (main.c:48)
==18066== 
==18066== Conditional jump or move depends on uninitialised value(s)
==18066==    at 0x48CF8D4: __printf_fp_buffer_1.isra.0 (printf_fp.c:794)
==18066==    by 0x48D135B: __printf_fp_l_buffer (printf_fp.c:1122)
==18066==    by 0x48D86C4: __printf_fp_spec (vfprintf-internal.c:266)
==18066==    by 0x48D86C4: __printf_buffer (vfprintf-internal.c:999)
==18066==    by 0x48D973A: __vfprintf_internal (vfprintf-internal.c:1544)
==18066==    by 0x49A5C9C: __printf_chk (printf_chk.c:33)
==18066==    by 0x1093A6: main (main.c:59)
==18066==  Uninitialised value was created by a heap allocation
==18066==    at 0x4846828: malloc (in /usr/libexec/valgrind/vgpreload_memcheck-amd64-linux.so)
==18066==    by 0x1098F5: aplicar_mh (mh.c:72)
==18066==    by 0x109339: main (main.c:48)
==18066== 
==18066== Conditional jump or move depends on uninitialised value(s)
==18066==    at 0x48CFAFB: __printf_fp_buffer_1.isra.0 (printf_fp.c:805)
==18066==    by 0x48D135B: __printf_fp_l_buffer (printf_fp.c:1122)
==18066==    by 0x48D86C4: __printf_fp_spec (vfprintf-internal.c:266)
==18066==    by 0x48D86C4: __printf_buffer (vfprintf-internal.c:999)
==18066==    by 0x48D973A: __vfprintf_internal (vfprintf-internal.c:1544)
==18066==    by 0x49A5C9C: __printf_chk (printf_chk.c:33)
==18066==    by 0x1093A6: main (main.c:59)
==18066==  Uninitialised value was created by a heap allocation
==18066==    at 0x4846828: malloc (in /usr/libexec/valgrind/vgpreload_memcheck-amd64-linux.so)
==18066==    by 0x1098F5: aplicar_mh (mh.c:72)
==18066==    by 0x109339: main (main.c:48)
==18066== 
==18066== Conditional jump or move depends on uninitialised value(s)
==18066==    at 0x48CFB26: __printf_fp_buffer_1.isra.0 (printf_fp.c:808)
==18066==    by 0x48D135B: __printf_fp_l_buffer (printf_fp.c:1122)
==18066==    by 0x48D86C4: __printf_fp_spec (vfprintf-internal.c:266)
==18066==    by 0x48D86C4: __printf_buffer (vfprintf-internal.c:999)
==18066==    by 0x48D973A: __vfprintf_internal (vfprintf-internal.c:1544)
==18066==    by 0x49A5C9C: __printf_chk (printf_chk.c:33)
==18066==    by 0x1093A6: main (main.c:59)
==18066==  Uninitialised value was created by a heap allocation
==18066==    at 0x4846828: malloc (in /usr/libexec/valgrind/vgpreload_memcheck-amd64-linux.so)
==18066==    by 0x1098F5: aplicar_mh (mh.c:72)
==18066==    by 0x109339: main (main.c:48)
==18066== 
==18066== Conditional jump or move depends on uninitialised value(s)
==18066==    at 0x48CFB5C: __printf_fp_buffer_1.isra.0 (printf_fp.c:818)
==18066==    by 0x48D135B: __printf_fp_l_buffer (printf_fp.c:1122)
==18066==    by 0x48D86C4: __printf_fp_spec (vfprintf-internal.c:266)
==18066==    by 0x48D86C4: __printf_buffer (vfprintf-internal.c:999)
==18066==    by 0x48D973A: __vfprintf_internal (vfprintf-internal.c:1544)
==18066==    by 0x49A5C9C: __printf_chk (printf_chk.c:33)
==18066==    by 0x1093A6: main (main.c:59)
==18066==  Uninitialised value was created by a heap allocation
==18066==    at 0x4846828: malloc (in /usr/libexec/valgrind/vgpreload_memcheck-amd64-linux.so)
==18066==    by 0x1098F5: aplicar_mh (mh.c:72)
==18066==    by 0x109339: main (main.c:48)
==18066== 
==18066== Conditional jump or move depends on uninitialised value(s)
==18066==    at 0x48CFF99: round_away (rounding-mode.h:52)
==18066==    by 0x48CFF99: __printf_fp_buffer_1.isra.0 (printf_fp.c:825)
==18066==    by 0x48D135B: __printf_fp_l_buffer (printf_fp.c:1122)
==18066==    by 0x48D86C4: __printf_fp_spec (vfprintf-internal.c:266)
==18066==    by 0x48D86C4: __printf_buffer (vfprintf-internal.c:999)
==18066==    by 0x48D973A: __vfprintf_internal (vfprintf-internal.c:1544)
==18066==    by 0x49A5C9C: __printf_chk (printf_chk.c:33)
==18066==    by 0x1093A6: main (main.c:59)
==18066==  Uninitialised value was created by a heap allocation
==18066==    at 0x4846828: malloc (in /usr/libexec/valgrind/vgpreload_memcheck-amd64-linux.so)
==18066==    by 0x1098F5: aplicar_mh (mh.c:72)
==18066==    by 0x109339: main (main.c:48)
==18066== 
==18066== Conditional jump or move depends on uninitialised value(s)
==18066==    at 0x48CFCDD: __printf_fp_buffer_1.isra.0 (printf_fp.c:987)
==18066==    by 0x48D135B: __printf_fp_l_buffer (printf_fp.c:1122)
==18066==    by 0x48D86C4: __printf_fp_spec (vfprintf-internal.c:266)
==18066==    by 0x48D86C4: __printf_buffer (vfprintf-internal.c:999)
==18066==    by 0x48D973A: __vfprintf_internal (vfprintf-internal.c:1544)
==18066==    by 0x49A5C9C: __printf_chk (printf_chk.c:33)
==18066==    by 0x1093A6: main (main.c:59)
==18066==  Uninitialised value was created by a heap allocation
==18066==    at 0x4846828: malloc (in /usr/libexec/valgrind/vgpreload_memcheck-amd64-linux.so)
==18066==    by 0x1098F5: aplicar_mh (mh.c:72)
==18066==    by 0x109339: main (main.c:48)
==18066== 
==18066== Conditional jump or move depends on uninitialised value(s)
==18066==    at 0x48CFD33: __printf_fp_buffer_1.isra.0 (printf_fp.c:999)
==18066==    by 0x48D135B: __printf_fp_l_buffer (printf_fp.c:1122)
==18066==    by 0x48D86C4: __printf_fp_spec (vfprintf-internal.c:266)
==18066==    by 0x48D86C4: __printf_buffer (vfprintf-internal.c:999)
==18066==    by 0x48D973A: __vfprintf_internal (vfprintf-internal.c:1544)
==18066==    by 0x49A5C9C: __printf_chk (printf_chk.c:33)
==18066==    by 0x1093A6: main (main.c:59)
==18066==  Uninitialised value was created by a heap allocation
==18066==    at 0x4846828: malloc (in /usr/libexec/valgrind/vgpreload_memcheck-amd64-linux.so)
==18066==    by 0x1098F5: aplicar_mh (mh.c:72)
==18066==    by 0x109339: main (main.c:48)
==18066== 
==18066== Syscall param write(buf) points to uninitialised byte(s)
==18066==    at 0x498A844: write (write.c:26)
==18066==    by 0x4901A54: _IO_file_write@@GLIBC_2.2.5 (fileops.c:1181)
==18066==    by 0x4900650: new_do_write (fileops.c:449)
==18066==    by 0x4900650: _IO_do_write@@GLIBC_2.2.5 (fileops.c:426)
==18066==    by 0x49042C5: _IO_flush_all (genops.c:705)
==18066==    by 0x49048EC: _IO_cleanup (genops.c:843)
==18066==    by 0x48B5B73: __run_exit_handlers (exit.c:129)
==18066==    by 0x48B5BBD: exit (exit.c:138)
==18066==    by 0x48981D0: (below main) (libc_start_call_main.h:74)
==18066==  Address 0x4c98d60 is 1,616 bytes inside a block of size 4,096 alloc'd
==18066==    at 0x4846828: malloc (in /usr/libexec/valgrind/vgpreload_memcheck-amd64-linux.so)
==18066==    by 0x48F3294: _IO_file_doallocate (filedoalloc.c:101)
==18066==    by 0x4903603: _IO_doallocbuf (genops.c:347)
==18066==    by 0x490106F: _IO_file_overflow@@GLIBC_2.2.5 (fileops.c:745)
==18066==    by 0x4901B8E: _IO_new_file_xsputn (fileops.c:1244)
==18066==    by 0x4901B8E: _IO_file_xsputn@@GLIBC_2.2.5 (fileops.c:1197)
==18066==    by 0x48CECC8: __printf_buffer_flush_to_file (printf_buffer_to_file.c:59)
==18066==    by 0x48CECC8: __printf_buffer_to_file_done (printf_buffer_to_file.c:120)
==18066==    by 0x48D9742: __vfprintf_internal (vfprintf-internal.c:1545)
==18066==    by 0x49A5C9C: __printf_chk (printf_chk.c:33)
==18066==    by 0x10938B: printf (stdio2.h:86)
==18066==    by 0x10938B: main (main.c:52)
==18066==  Uninitialised value was created by a heap allocation
==18066==    at 0x4846828: malloc (in /usr/libexec/valgrind/vgpreload_memcheck-amd64-linux.so)
==18066==    by 0x1098F5: aplicar_mh (mh.c:72)
==18066==    by 0x109339: main (main.c:48)
==18066== 
==18066== 
==18066== HEAP SUMMARY:
==18066==     in use at exit: 2,159,600 bytes in 201 blocks
==18066==   total heap usage: 206 allocs, 5 frees, 2,171,792 bytes allocated
==18066== 
==18066== 161,600 (1,600 direct, 160,000 indirect) bytes in 100 blocks are definitely lost in loss record 2 of 3
==18066==    at 0x4846828: malloc (in /usr/libexec/valgrind/vgpreload_memcheck-amd64-linux.so)
==18066==    by 0x1098F5: aplicar_mh (mh.c:72)
==18066==    by 0x109339: main (main.c:48)
==18066== 
==18066== 1,998,000 bytes in 1 blocks are definitely lost in loss record 3 of 3
==18066==    at 0x4846828: malloc (in /usr/libexec/valgrind/vgpreload_memcheck-amd64-linux.so)
==18066==    by 0x109561: read_distances (io.c:14)
==18066==    by 0x1092DE: main (main.c:34)
==18066== 
==18066== LEAK SUMMARY:
==18066==    definitely lost: 1,999,600 bytes in 101 blocks
==18066==    indirectly lost: 160,000 bytes in 100 blocks
==18066==      possibly lost: 0 bytes in 0 blocks
==18066==    still reachable: 0 bytes in 0 blocks
==18066==         suppressed: 0 bytes in 0 blocks
==18066== 
==18066== For lists of detected and suppressed errors, rerun with: -s
==18066== ERROR SUMMARY: 63941 errors from 34 contexts (suppressed: 0 from 0)

    \end{verbatim}



\begin{thebibliography}{9}

\bibitem{baeldung_parallel}
Baeldung, \emph{A Quick Guide to the Linux parallel Command}, 2025. [En línea]. Disponible en: \url{https://www.baeldung.com/linux/parallel-command-tutorial} [Accedido: 09-sep-2026].

\bibitem{lemon_tendencias}
Lemon Learning, \emph{10 tendencias IT}, s.f. [En línea]. Disponible en: \url{https://lemonlearning.com/es/blog/10-tendencias-it} [Accedido: 09-sep-2026].

\end{thebibliography}










\end{document}